\scanpage{52}
homéomorphismes universels\ldots

Les réflexions précédentes suggèrent aussi des énoncés comme \unsure{les|le} \unsure{suivants|suivant} : Pour un schém. de base $S$ loc. noeth. donné,
\margin{$S$ de car. $0$ ?}
le foncteur $X\unsure{-|\longrightarrow}X_{\mathrm{et}}$, allant de la catégorie des schémas\inserted{ réduits} loc. de prés. finie \inserted{\deleted{réduits}} sur $S$, vers la 2-catégorie des topos au dessus de $X_{\mathrm{et}}$, est \deleted{pleinement} 1-fidèle (deux hom $f,g:X=Y$ tels que les homomorphismes de topos \unsure{\inserted{$f_{\mathrm{et}},g_{\mathrm{et}}$}} $X_{\mathrm{et}}\rightrightarrows Y_{\mathrm{et}}$ au-dessus de $S_{\mathrm{et}}$ soient \unsure{ég|$fg$|\deleted{ég}|\deleted{$fg$}} isomorphes, sont égaux) et $=$ peut-être \uline{pleinement} \uline{fidèle}, quand on passe à la catégorie des fractions, de $(\mathrm{Sch}_{\mathrm{l.t.f}})/S$ obtenue en rendant inv. les homéom. universels\ldots \unsure{Exprimant|Expérimenter} \deleted{\unsure{p}} ceci p.ex. pour les autom d'une courbe alg. propre sur une ext. finie de $\mathbb{Q}$, \unsure{en|on} \inserted{re}\unsure{trouverait|trouvant|trouvait} le ``fait'' que tout autom ext. de $\mathcal E_K$ ($K$ le corps des fonctions de $X$) qui respecte la structure à lacets et qui commute à l'action de $\Gamma$, provient d'un autom de $X$
\margin{51}
